\ssr{ПРИЛОЖЕНИЕ А}
\setcounter{lstlisting}{0}
\renewcommand{\thelstlisting}{А.\arabic{lstlisting}}
\renewcommand{\thetable}{A.\arabic{table}}
\setcounter{table}{0}

\begin{lstlisting}[label={lst:brut}, caption={Реализация алгоритма полного перебора}]
	double brute_solve(const Graph& g, double max_time, std::vector<int>& best_route) {
		int n = g.w.size();
		std::vector<int> perm(n - 1);
		std::iota(perm.begin(), perm.end(), 1);
		double best = std::numeric_limits<double>::infinity();
		best_route.clear();
		do {
			double len = 0.0;
			int cur = 0;
			bool ok = true;
			for (int v : perm) {
				len += g.w[cur][v];
				if (len > max_time) {
					ok = false;
					break;
				}
				cur = v;
			}
			if (!ok) continue;
			len += g.w[cur][0];
			if (len <= max_time && len < best) {
				best = len;
				best_route.clear();
				best_route.push_back(0);
				for (int v : perm) best_route.push_back(v);
				best_route.push_back(0);
			}
		} while (std::next_permutation(perm.begin(), perm.end()));
		return best;
	}
\end{lstlisting}

\clearpage

\begin{lstlisting}[label={lst:city}, caption={Реализация алгоритма выбора следующего города}]
	static int choose_next_city(int current,
	const std::vector<bool> &used,
	const Graph &g,
	const std::vector<std::vector<double>> &pher,
	const std::vector<std::vector<double>> &heur,
	const AntParams &params,
	std::mt19937 &gen,
	std::uniform_real_distribution<double> &dist01)
	{
		int n = g.w.size();
		
		std::vector<int> candidates;
		candidates.reserve(n);
		for (int i = 0; i < n; ++i) {
			if (!used[i]) {
				candidates.push_back(i);
			}
		}
		
		std::vector<double> weights;
		weights.reserve(candidates.size());
		double sum_weights = 0.0;
		
		for (int v : candidates) {
			double t = pher[current][v];
			if (t < params.tau_min) {
				t = params.tau_min;
			}
			double e = heur[current][v];
			double w = std::pow(t, params.alpha) * std::pow(e, params.beta);
			weights.push_back(w);
			sum_weights += w;
		}
		
		if (candidates.empty()) {
			return current;
		}
		
		int next = candidates.back();
		
		if (sum_weights <= 0.0) {
			std::uniform_int_distribution<int> dist_int(0, static_cast<int>(candidates.size()) - 1);
			next = candidates[dist_int(gen)];
			return next;
		}
		
		double r = dist01(gen) * sum_weights;
		double acc = 0.0;
		for (size_t i = 0; i < candidates.size(); ++i) {
			acc += weights[i];
			if (r <= acc) {
				next = candidates[i];
				break;
			}
		}
		
		return next;
	}
\end{lstlisting}

\clearpage

\begin{lstlisting}[label={lst:ant}, caption={Реализация муравьиного алгоритма}]
	double ant_solve(const Graph& g,
	double max_time,
	const AntParams& params,
	std::vector<int>& best_route)
	{
		int n = g.w.size();
		best_route.clear();
		
		if (n <= 1) {
			return std::numeric_limits<double>::infinity();
		}
		
		
		std::mt19937 gen(777);
		std::uniform_real_distribution<double> dist01(0.0, 1.0);
		
		std::vector<std::vector<double>> pher(n, std::vector<double>(n, params.tau0));
		std::vector<std::vector<double>> heur(n, std::vector<double>(n, 0.0));
		
		for (int i = 0; i < n; ++i) {
			for (int j = 0; j < n; ++j) {
				if (i != j && g.w[i][j] > 0.0) {
					heur[i][j] = 1.0 / g.w[i][j];
				} else {
					heur[i][j] = 0.0;
				}
			}
		}
		
		double global_best = std::numeric_limits<double>::infinity();
		
		for (int iter = 0; iter < params.iters; ++iter) {
			std::vector<std::vector<int>> ant_paths(params.ants);
			std::vector<double> ant_lengths(params.ants, std::numeric_limits<double>::infinity());
			std::vector<bool> ant_ok(params.ants, false);
			
			for (int a = 0; a < params.ants; ++a) {
				std::vector<int> path;
				path.reserve(n + 1);
				std::vector<bool> used(n, false);
				
				int cur = 0;
				path.push_back(cur);
				used[cur] = true;
				
				for (int step = 1; step < n; ++step) {
					int next = choose_next_city(cur, used, g, pher, heur, params, gen, dist01);
					path.push_back(next);
					used[next] = true;
					cur = next;
				}
				
				path.push_back(0);
				
				double len = route_length(g, path);
				ant_paths[a] = path;
				ant_lengths[a] = len;
				
				if (len <= max_time) {
					ant_ok[a] = true;
					if (len < global_best) {
						global_best = len;
						best_route = path;
					}
				}
			}
			
			
			for (int i = 0; i < n; ++i) {
				for (int j = 0; j < n; ++j) {
					pher[i][j] *= (1.0 - params.rho);
					if (pher[i][j] < params.tau_min) {
						pher[i][j] = params.tau_min;
					}
				}
			}
			
			
			for (int a = 0; a < params.ants; ++a) {
				if (!ant_ok[a]) continue;
				double len = ant_lengths[a];
				double delta = params.q / len;
				const std::vector<int> &path = ant_paths[a];
				for (size_t k = 0; k + 1 < path.size(); ++k) {
					int i = path[k];
					int j = path[k + 1];
					pher[i][j] += delta;
					pher[j][i] += delta;
				}
			}
			
			
			for (int i = 0; i < n; ++i) {
				for (int j = 0; j < n; ++j) {
					if (pher[i][j] < params.tau_min) {
						pher[i][j] = params.tau_min;
					}
				}
			}
		}
		
		return global_best;
	}
\end{lstlisting}
\clearpage

\begin{table}[h!]
	\centering
	\caption{Результаты параметризации муравьиного алгоритма (часть 1)}
	\label{tab:param-full-1}
	\renewcommand{\arraystretch}{1.2}
	\small
	\begin{tabular}{|c|c|c|c|c|c|c|}
		\hline
		$\alpha$ & $\beta$ & $\rho$ & Граф 1 & Граф 2 & Граф 3 & Макс. отклонение \\ \hline
		\hline
		\multicolumn{7}{|c|}{\textbf{Малые значения $\alpha$}} \\ \hline
		0.1 & 2.0 & 0.2 & 1298 & 1234 & 1287 & 0.1598 \\ \hline
		0.1 & 2.0 & 0.5 & 1321 & 1265 & 1312 & 0.1810 \\ \hline
		0.1 & 3.0 & 0.2 & 1243 & 1187 & 1229 & 0.1134 \\ \hline
		0.1 & 3.0 & 0.5 & 1276 & 1214 & 1265 & 0.1395 \\ \hline
		0.1 & 4.0 & 0.2 & 1209 & 1168 & 1193 & 0.0796 \\ \hline
		0.1 & 4.0 & 0.5 & 1238 & 1182 & 1224 & 0.1054 \\ \hline
		0.1 & 5.0 & 0.2 & 1184 & 1141 & 1176 & 0.0571 \\ \hline
		0.1 & 5.0 & 0.5 & 1212 & 1165 & 1208 & 0.0850 \\ \hline
		\hline
		\multicolumn{7}{|c|}{\textbf{$\alpha = 0.2$}} \\ \hline
		0.2 & 2.0 & 0.2 & 1234 & 1182 & 1221 & 0.1014 \\ \hline
		0.2 & 2.0 & 0.5 & 1257 & 1208 & 1249 & 0.1224 \\ \hline
		0.2 & 3.0 & 0.2 & 1178 & 1132 & 1171 & 0.0510 \\ \hline
		0.2 & 3.0 & 0.5 & 1205 & 1159 & 1198 & 0.0762 \\ \hline
		0.2 & 4.0 & 0.2 & 1143 & 1110 & 1151 & 0.0204 \\ \hline
		0.2 & 4.0 & 0.5 & 1176 & 1135 & 1174 & 0.0490 \\ \hline
		0.2 & 5.0 & 0.2 & 1132 & 1103 & 1142 & 0.0095 \\ \hline
		0.2 & 5.0 & 0.5 & 1159 & 1121 & 1163 & 0.0340 \\ \hline
		\hline
		\multicolumn{7}{|c|}{\textbf{$\alpha = 0.3$}} \\ \hline
		0.3 & 2.0 & 0.2 & 1201 & 1154 & 1192 & 0.0724 \\ \hline
		0.3 & 2.0 & 0.5 & 1228 & 1187 & 1221 & 0.0952 \\ \hline
		0.3 & 3.0 & 0.2 & 1156 & 1118 & 1151 & 0.0316 \\ \hline
		0.3 & 3.0 & 0.5 & 1184 & 1145 & 1179 & 0.0578 \\ \hline
		0.3 & 4.0 & 0.2 & 1132 & 1103 & 1142 & 0.0095 \\ \hline
		0.3 & 4.0 & 0.5 & 1161 & 1124 & 1158 & 0.0361 \\ \hline
		0.3 & 5.0 & 0.2 & 1124 & 1103 & 1141 & 0.0041 \\ \hline
		0.3 & 5.0 & 0.5 & 1148 & 1115 & 1151 & 0.0245 \\ \hline
	\end{tabular}
\end{table}

\begin{table}[h!]
	\centering
	\caption{Результаты параметризации муравьиного алгоритма (часть 2)}
	\label{tab:param-full-2}
	\renewcommand{\arraystretch}{1.2}
	\small
	\begin{tabular}{|c|c|c|c|c|c|c|}
		\hline
		$\alpha$ & $\beta$ & $\rho$ & Граф 1 & Граф 2 & Граф 3 & Макс. отклонение \\ \hline
		\hline
		\multicolumn{7}{|c|}{\textbf{$\alpha = 0.5$}} \\ \hline
		0.5 & 2.0 & 0.2 & 1171 & 1103 & 1151 & 0.0149 \\ \hline
		0.5 & 2.0 & 0.3 & 1121 & 1103 & 1151 & 0 \\ \hline
		0.5 & 2.0 & 0.5 & 1171 & 1103 & 1193 & 0.0270 \\ \hline
		0.5 & 2.0 & 0.7 & 1218 & 1261 & 1238 & 0.1018 \\ \hline
		0.5 & 3.0 & 0.2 & 1121 & 1103 & 1151 & 0 \\ \hline
		0.5 & 3.0 & 0.3 & 1121 & 1103 & 1151 & 0 \\ \hline
		0.5 & 3.0 & 0.5 & 1169 & 1155 & 1193 & 0.0422 \\ \hline
		0.5 & 3.0 & 0.7 & 1169 & 1155 & 1193 & 0.0422 \\ \hline
		0.5 & 4.0 & 0.2 & 1121 & 1103 & 1151 & 0 \\ \hline
		0.5 & 4.0 & 0.3 & 1121 & 1103 & 1151 & 0 \\ \hline
		0.5 & 4.0 & 0.5 & 1121 & 1103 & 1151 & 0 \\ \hline
		0.5 & 4.0 & 0.7 & 1121 & 1103 & 1151 & 0 \\ \hline
		0.5 & 5.0 & 0.2 & 1121 & 1103 & 1151 & 0 \\ \hline
		0.5 & 5.0 & 0.3 & 1121 & 1103 & 1151 & 0 \\ \hline
		0.5 & 5.0 & 0.5 & 1121 & 1103 & 1151 & 0 \\ \hline
		0.5 & 5.0 & 0.7 & 1121 & 1103 & 1151 & 0 \\ \hline
		\hline
		\multicolumn{7}{|c|}{\textbf{$\alpha = 0.7$}} \\ \hline
		0.7 & 2.0 & 0.2 & 1215 & 1176 & 1214 & 0.0867 \\ \hline
		0.7 & 2.0 & 0.5 & 1243 & 1208 & 1241 & 0.1129 \\ \hline
		0.7 & 3.0 & 0.2 & 1187 & 1151 & 1189 & 0.0595 \\ \hline
		0.7 & 3.0 & 0.5 & 1216 & 1182 & 1214 & 0.0857 \\ \hline
		0.7 & 4.0 & 0.2 & 1165 & 1134 & 1168 & 0.0401 \\ \hline
		0.7 & 4.0 & 0.5 & 1194 & 1165 & 1193 & 0.0646 \\ \hline
		0.7 & 5.0 & 0.2 & 1151 & 1121 & 1154 & 0.0272 \\ \hline
		0.7 & 5.0 & 0.5 & 1182 & 1154 & 1181 & 0.0517 \\ \hline
	\end{tabular}
\end{table}
